\chapter{Testbare Vorhersagen}

\section{Einleitung}

Die WDBT+ macht eine Reihe spezifischer, von der Allgemeinen Relativitätstheorie und dem Standardmodell der Kosmologie abweichender Vorhersagen. Diese Vorhersagen sind prinzipiell experimentell überprüfbar – entweder mit bestehenden oder mit zukünftigen Technologien.

Die folgende Liste fasst die wichtigsten testbaren Vorhersagen zusammen, geordnet nach physikalischen Bereichen.

\section{Kosmologie und Galaxien}

\subsection{Fraktale Skalengesetze}

Die fraktale Raumdimension \( D \approx 2,704 \) führt zu gebrochenen Exponenten in kosmologischen Skalengesetzen:

\begin{itemize}
    \item Galaxien-Dichteprofil: \( \rho(r) \propto r^{-(3-D)} \approx r^{-0,296} \) (flacher als das NFW-Profil)
    \item Rotationskurven: \( v(r) \propto r^{(D-2)/2} \approx r^{0,352} \) (bei sehr großen Radien steigt die Kurve langsam an, fällt nicht ab)
    \item CMB-Korrelationen: \( C(\theta) \propto \theta^{-(3-D)} \approx \theta^{-0,296} \) mit einer festen, aus der Dodekaeder-Symmetrie folgenden Achse („Achse des Bösen“)
    \item Materieentstehung: \( \frac{dM}{d\ln r} \propto r^{3-D} \approx r^{0,296} \)
\end{itemize}

\subsection{Maximale Masse kompakter Objekte}

Die Theorie sagt eine natürliche Obergrenze für kompakte Objekte voraus:

\[
M_{\text{BEC}} \approx 3 \times 10^6 M_{\odot}
\]

Beobachtete supermassereiche Objekte mit Massen über dieser Grenze sind keine einzelnen kompakten Objekte, sondern Systeme aus einem BEC-Kern und einer massiven Umgebung.

\subsection{Doppelkern-Systeme}

Bei Galaxienverschmelzungen sagt die Theorie voraus, dass Doppelkern-Systeme häufiger sind als in der ART vorhergesagt, da BEC-Objekte bei Überschreiten der maximalen Masse zerreißen können, anstatt zu verschmelzen.

\subsection{Hubble-Gesetz}

Die Rotverschiebung ist keine Doppler-Folge der Expansion, sondern eine kumulative gravitative Wirkung. Daher ergeben sich Abweichungen vom linearen Hubble-Gesetz auf großen Skalen.

\subsection{Übersicht der testbaren Vorhersagen}

\subsection{Übersicht der testbaren Vorhersagen}

Die folgende Tabelle fasst die wichtigsten Vorhersagen zusammen:

\begin{tabular}{|l|l|l|}
\hline
\textbf{Vorhersage} & \textbf{Abweichung von ART} & \textbf{Experiment} \\
\hline
Lichtablenkung & \(\Delta\phi = \Delta\phi_{\text{ART}}(1 + \alpha\lambda^2)\) & EHT, VLBI, LISA \\
Gravitationswellen-Dispersion & \(v_{\text{phase}} = c \left(1 + \beta \left(\frac{f}{f_0}\right)^{-0,296}\right)\) & LISA, Einstein-Teleskop \\
Neutrinomasse & \(m_\nu \approx 0,1\,\text{eV}/c^2\) & KATRIN, JUNO, DUNE \\
\textbf{Hubble-Skalenabhängigkeit} & \(\mathbf{H_{\text{eff}}(d) = H_0 (d/d_0)^{0,704}}\) & Supernovae, Quasare, GW \\
Maximale kompakte Masse & \(M_{\text{BEC}} \approx 3 \times 10^6 M_\odot\) & Galaxienzentren, EHT \\
Galaxien-Dichteprofil & \(\rho(r) \propto r^{-0,296}\) & EUCLID, Rubin (LSST) \\
Rotationskurven & \(v(r) \propto r^{0,352}\) (asymptotisch steigend) & Extrem breite Galaxien \\
\hline
\end{tabular}

Die benötigte relative Genauigkeit liegt für die meisten Experimente im Bereich \( 10^{-6} \) bis \( 10^{-2} \). Jede dieser Vorhersagen ist prinzipiell falsifizierbar.

\section{Gravitationsphysik}

\subsection{Frequenzabhängige Lichtablenkung}

Die ART sagt eine frequenzunabhängige Lichtablenkung voraus. Die WDBT+ sagt dagegen eine frequenzabhängige Ablenkung voraus:

\[
\Delta\phi(\nu) = \Delta\phi_0 \left(1 + \alpha \frac{\nu_0}{\nu}\right) \propto \lambda^2
\]

Dabei ist \( \alpha \) eine dimensionslose Konstante, die aus der fraktalen Raumdimension \( D \) folgt (siehe Anhang D, Gleichung D.12). Dieser Effekt ist mit hochpräzisen Interferometern (z. B. LISA) oder spektral aufgelösten Messungen am Sonnenrand testbar.

\subsection{Gravitationswellen-Dispersion}

Die fraktale Raumstruktur führt zu einer frequenzabhängigen Ausbreitungsgeschwindigkeit von Gravitationswellen:

\[
v_{\text{phase}}(\omega) = c \left(1 + \beta \left(\frac{\omega}{\omega_0}\right)^{D-3} + \dots\right)
\]

Dabei ist \( \beta \) eine dimensionslose Konstante, die aus der fraktalen Raumdimension \( D \) folgt (siehe Anhang I, Gleichung I.16). Mit \( D \approx 2,704 \) ist \( D-3 \approx -0,296 \), also eine schwache Dispersion bei hohen Frequenzen. Dies ist mit zukünftigen Detektoren wie LISA oder dem Einstein-Teleskop testbar.

\subsection{Shapiro-Laufzeitverzögerung}

Die ART sagt eine frequenzunabhängige Laufzeitverzögerung voraus. Die WDBT+ sagt eine zusätzliche Frequenzabhängigkeit voraus:

\[
\Delta t_{\text{WG}} \propto \lambda^{-2}
\]

Dies ist mit Pulsar-Timing-Experimenten testbar.

\subsection{Spiralbahnen statt Ellipsen}

Die DSTT sagt voraus, dass alle Bahnen logarithmische Spiralen sind, keine geschlossenen Ellipsen. Die radiale Drift ist extrem langsam:

\[
\frac{\dot{r}}{r} \approx 10^{-11} \text{ pro Jahr}
\]

Für den Mond ergibt sich eine Exzentrizitätsabnahme von \( \dot{e} \approx -2,7 \times 10^{-11} \) pro Jahr. Dieser Effekt ist mit heutigen Methoden (Lunar Laser Ranging) nicht nachweisbar, aber prinzipiell testbar.

\section{Quanten- und Teilchenphysik}

\subsection{Neutrinomasse}

Die Theorie sagt eine Neutrinomasse von

\[
m_\nu \approx 0,1 \, \text{eV}/c^2
\]

voraus, direkt aus der fundamentalen Länge \( l_0 \). Dies ist mit Experimenten wie KATRIN oder dem neutrinolosen Doppelbeta-Zerfall testbar.

\subsection{Neutrino-Oszillationen}

Die fraktale Dispersionsrelation führt zu einer charakteristischen Skalierung der Oszillationslänge:

\[
L_{\text{osc}} \propto E^{0,775}
\]

Dies weicht von den Vorhersagen des Standardmodells ab und ist mit Neutrino-Oszillationsexperimenten (JUNO, DUNE, Hyper-Kamiokande) testbar.

\subsection{Drei Neutrino-Generationen}

Die drei Generationen sind Moden des Q-Feldes im fraktalen Raum. Die Massenunterschiede zwischen den Generationen folgen aus einer feinen Aufspaltung durch die fraktale Geometrie.

\subsection{Modifizierter Lamb-Shift}

Die Wechselwirkung des Elektrons mit dem Quantenpotential des Vakuums führt zu einem Zusatzterm im Lamb-Shift:

\[
\Delta E_{\text{Lamb}}^{\text{WDBT}} = \Delta E_{\text{QED}} + \frac{e^2 \hbar}{4 \pi \epsilon_0 m_e^2 c^3} \langle r \rangle
\]

Dabei ist \( \langle r \rangle \) der Erwartungswert des Ortes des Elektrons im gebundenen Zustand. Dieser Term ist klein, aber prinzipiell messbar.

\section{Plasmaphysik}

\subsection{Fraktale Skalierung im Sonnenwind}

Fluktuationen des Sonnenwinds sollten einen Exponenten \( 3-D \approx 0,296 \) zeigen – testbar mit bestehenden Satellitendaten.

\subsection{Sonnenmasse}

Die Theorie sagt voraus, dass die Sonne keinen messbaren Masseverlust durch den Sonnenwind erfährt, da der Sonnenwind aus neu entstehender Materie besteht, nicht aus der Sonne selbst.

\subsection{Stromdichten in Plasmen}

In Laborplasmen und astrophysikalischen Plasmen sollten Stromdichten mit \( j(r) \propto r^{D-3} \approx r^{-0,296} \) skalieren.

\section{Zusammenfassung der wichtigsten Vorhersagen}

Die folgende Tabelle fasst die wichtigsten testbaren Vorhersagen zusammen:

\[
\begin{array}{|l|l|l|}
\hline
\text{Vorhersage} & \text{Wert / Skalierung} & \text{Testmethode} \\
\hline
\text{Galaxien-Dichteprofil} & \rho(r) \propto r^{-0,296} & \text{EUCLID, Rubin (LSST)} \\
\text{Rotationskurven} & v(r) \propto r^{0,352} & \text{Extrem breite Galaxien} \\
\text{CMB-Korrelationen} & C(\theta) \propto \theta^{-0,296} & \text{Planck, LiteBIRD} \\
\text{Neutrinomasse} & m_\nu \approx 0,1 \, \text{eV}/c^2 & \text{KATRIN, Doppelbeta-Zerfall} \\
\text{Sonnenmasse} & \text{Kein Masseverlust durch Wind} & \text{Präzise Sonnenbeobachtung} \\
\text{Gravitationswellen-Dispersion} & v_{\text{phase}}(\omega) = c(1 + \alpha \omega^{-0,296}) & \text{LISA, Einstein-Teleskop} \\
\text{Skalierung im Plasma} & j(r) \propto r^{-0,296} & \text{Laborplasmen, Sonnenwind} \\
\text{Lichtablenkung} & \Delta\phi \propto \lambda^2 & \text{Multiband-Beobachtungen} \\
\text{Shapiro-Laufzeit} & \text{Frequenzabhängig} & \text{Pulsar-Timing} \\
\text{Maximale Kernmasse} & M_{\text{BEC}} \approx 3 \times 10^6 M_\odot & \text{Beobachtung von Galaxienzentren} \\
\text{EHT-Schatten} & r_{\text{ph}} = 3GM/c^2, \text{ andere Feinstruktur} & \text{EHT, nächste Generation} \\
\hline
\end{array}
\]

\section{Fazit}

Die WDBT+ macht spezifische, testbare Vorhersagen, die sie von der Allgemeinen Relativitätstheorie und dem Standardmodell der Kosmologie unterscheiden. Sollten diese Vorhersagen bestätigt werden, wäre dies ein starkes Indiz für die Richtigkeit der Theorie. Sollten sie widerlegt werden, wäre die Theorie falsifiziert – ein Zeichen ihrer wissenschaftlichen Qualität.